% Shared abstract (150-250 words; no citations, footnotes or abbreviations).
A compiler that must run inside a fixed memory budget, for example on a microcontroller, needs a symbol table whose size is bounded in advance, whose behavior at the bound is defined, and whose cost per call cannot be driven up by hostile source text. Scoped symbol tables in textbooks and in production compilers rely on a general-purpose allocator and offer none of these properties. This article presents STRATA, a scoped symbol table that keeps its index, its interned names, and an undo stack in one caller-supplied buffer. An operation that does not fit is refused and leaves the table bit-for-bit unchanged. We show that the state of the table after any sequence of operations equals that of a fresh table holding only the live declarations, and we use this property to bound the probability that a limit on index probes refuses a legitimate name. On symbol-table events logged by an instrumented C compiler for eight open-source projects, STRATA needs a fixed buffer \SaveLo\ to \SaveHi\ percent smaller than a conventional single-buffer table, with replay time within \ArenaAbsMax\ percent of it and no established difference on most traces, and \HeapMetisLo\ to \HeapMetisHi\ percent less heap than its multi-allocation predecessor. Identifiers that collide in the default hash for every seed slow declarations by a factor of \FloodRatio; a keyed hash and a probe limit remove the slowdown. The guarantees are supported by differential fuzzing, exhaustive enumeration of short operation sequences, and a run on an emulated microcontroller.
